\vspace{2.5pt plus 3pt minus 0.5pt}\noindent\hspace*{1.2em}\kbvnum{18.}\hspace{0.55em}ಎಂಟು ವಿಧವಾದ ಅಕ್ಷರ ಗಣಗಳನ್ನು ಮತ್ತು ಅವುಗಳ ವಿನ್ಯಾಸವನ್ನು ಬರೆದು ತಿಳಿಸಿ\par

\vspace{2.5pt plus 3pt minus 0.5pt}\noindent\hspace*{1.2em}\kbvnum{19.}\hspace{0.55em}ಕೆಳಗೆ ಕೊಟ್ಟಿರುವ ಪದಗಳಿಗೆ ಪ್ರಸ್ತಾರ ಹಾಕಿ ಅವುಗಳು ಯಾವ ಅಕ್ಷರ ಗಣಕ್ಕೆ ಸೇರಿದವುಗಳೆಂಬುದನ್ನು ತಿಳಿಸಿರಿ.\par

\kbvsechead{ಕೆಂಜೇಳು ಸಾಧನೆ, ವಿನಾಶ ಮಕರ, ರೇಣುಕಾ, ಜನತಾ :}

\indent ಮಹಂತೇಶ್‌ ಶ್ರೀಲಂಕಾ, ಮಾಣಿಕ್ಯಂ\par

\vspace{2.5pt plus 3pt minus 0.5pt}\noindent\hspace*{1.2em}\kbvnum{20.}\hspace{0.55em}ಉತ್ಪಲಮಾಲಾ ವೃತ್ತ ಜಾತಿಯ ಪದ್ಯದ ಲಕ್ಷಣಗಳನ್ನು ತಿಳಿಸಿ.\par

\vspace{2.5pt plus 3pt minus 0.5pt}\noindent\hspace*{1.2em}\kbvnum{21.}\hspace{0.55em}ನಿಂದಿಸುವರ್‌ ದುರ್ಜನರ್‌ ಎಂ ।\par

\begin{kbvverse}
\noindent ಬೊಂದು ಭಯಂಜಿತ್ತು ಸುಕವಿ ರಚಿಸನೆ ಕೃತಿಯಂ ॥ \\
\noindent ಮಂದೇಹರ ಭಯದಿಂದರ । \\
\noindent ವಿಂದಸಖಂ ನಿಜ ಮಯೂಖಮಂ ಪ್ರಸರಿಸದೇ ? ॥
\end{kbvverse}

\noindent\hfill{\fontsize{9.5pt}{12.5pt}\selectfont\color{kbvgrey}— ರಾಜಶೇಖರ ವಿಲಾಸ}\par\vspace{3pt plus 2pt}

\vspace{2.5pt plus 3pt minus 0.5pt}\noindent\hspace*{1.2em}\kbvnum{೧)}\hspace{0.55em}ಈ ಪದ್ಯದ ತಾತ್ಪರ್ಯವೇನು ?\par

\vspace{2.5pt plus 3pt minus 0.5pt}\noindent\hspace*{1.2em}\kbvnum{೨)}\hspace{0.55em}ಈ ಪದ್ಯದಲ್ಲಿ ಯಾವ ಅಲಂಕಾರವಿರುವುದು ಎಂಬುದನ್ನು... ತಿಳಿಸಿರಿ,\par

\vspace{2.5pt plus 3pt minus 0.5pt}\noindent\hspace*{1.2em}\kbvnum{೩)}\hspace{0.55em}ಈ ಪದ್ಯಕ್ಕೆ ಪ್ರಸ್ತಾರ ಹಾಕಿ, ಮಾತ್ರಾ ಗಣಗಳನ್ನು ವಿಂಗಡಿಸಿರಿ,\par

\vspace{2.5pt plus 3pt minus 0.5pt}\noindent\hspace*{1.2em}\kbvnum{೪)}\hspace{0.55em}ಇದು “ಕಂದ” ಪದ್ಯ ಎಂಬುದನ್ನು ಹೇಗೆ ಸಮರ್ಥಿಸುವಿರಿ?\par
